
\documentclass[addpoints]{exam}
\usepackage{amsmath}
\usepackage{amssymb}
\usepackage{color}
\usepackage{booktabs}

% \printanswers

\newcommand{\event}{One-Variable Statistics}
\newcommand{\tournament}{}
\def\type{0}

\setlength\fillinlinelength{\textwidth}
\CorrectChoiceEmphasis{\color{red}\bfseries}
\SolutionEmphasis{\color{red}\bfseries}

\setlength\answerclearance{2pt}
\pagestyle{headandfoot}
\runningheadrule
\runningheader{\event}{\tournament}{\ifnum\type=1 Class Set Test \else Full Name:\hspace{1cm} \fi}
\runningfootrule
\runningfooter{}{Page \thepage\ of \numpages}{}

\begin{document}
\begin{center}
    {\huge Grade 12 Data Management \\ \event
    \ifprintanswers \space Key
    \else \space \ifnum\type<2 Test \fi \ifnum\type=1 \textbf{(CLASS SET)} \fi
          \ifnum\type=2 Answer Sheet \fi
    \fi \\ \vspace{3mm}\tournament\vspace{1mm}}
\end{center}

\vspace{.25\textheight}

\begin{center}
    First Name: \ifprintanswers
    \underline{\hspace{3.5cm}}\textbf{KEY}\underline{\hspace{5.5cm}}\\\vspace{5mm}
    \else \underline{\hspace{10cm}}\\ \vspace{5mm} \fi
    Last Name: \underline{\hspace{10cm}}\\ \vspace{5mm}
\end{center}

\noindent\textbf{\underline{Directions:}}
\begin{itemize}
    \ifnum\type=1 \item \textbf{This test is a class set.} Do not write on this paper. \fi
    \item Show all work for full marks. Round to 2 decimal places unless stated otherwise.
    \item The test is designed to be completed in 75 minutes.
\end{itemize}
\begin{center}
\textit{For grading use only}\\
\multirowgradetable{1}[pages]
\end{center}

\newpage
\begin{questions}

\section*{Part A --- Multiple Choice (15 marks)}
\vspace{5mm}

\question[1] Which measure of central tendency is most affected by extreme values (outliers)?
\begin{choices}
    \correctchoice Mean
    \choice Median
    \choice Mode
    \choice Range
\end{choices}

\question[1] The median of the data set $\{3, 7, 7, 8, 12, 15\}$ is:
\begin{choices}
    \choice $7$
    \correctchoice $7.5$
    \choice $8$
    \choice $8.67$
\end{choices}

\question[1] For a right-skewed distribution, the typical order of measures is:
\begin{choices}
    \choice mean $<$ median $<$ mode
    \correctchoice mode $<$ median $<$ mean
    \choice median $<$ mode $<$ mean
    \choice mean $=$ median $=$ mode
\end{choices}

\question[1] The first quartile $Q_1$ is:
\begin{choices}
    \choice The 25th data value
    \correctchoice The median of the lower half of the data
    \choice The average of the smallest and largest values
    \choice Equal to the 75th percentile
\end{choices}

\question[1] For the data $\{4, 6, 8, 10, 12, 14, 16, 18\}$, the interquartile range is:
\begin{choices}
    \choice $7$
    \choice $8$
    \correctchoice $8$
    \choice $14$
\end{choices}

\question[1] The population variance of $\{2, 4, 6, 8, 10\}$ is:
\begin{choices}
    \choice $4$
    \choice $6$
    \correctchoice $8$
    \choice $10$
\end{choices}

\question[1] Standard deviation is:
\begin{choices}
    \choice Always equal to the variance
    \correctchoice The square root of the variance
    \choice The square of the variance
    \choice Always greater than the mean
\end{choices}

\question[1] A data value is considered an outlier using the $1.5 \times \text{IQR}$ rule if it is:
\begin{choices}
    \choice More than 2 standard deviations from the mean
    \correctchoice Below $Q_1 - 1.5 \cdot \text{IQR}$ or above $Q_3 + 1.5 \cdot \text{IQR}$
    \choice Below $Q_1$ or above $Q_3$
    \choice More than 3 standard deviations from the mean
\end{choices}

\question[1] A student scores at the 80th percentile on a test. This means:
\begin{choices}
    \choice They scored 80\% on the test.
    \correctchoice 80\% of students scored at or below this score.
    \choice They answered 80 questions correctly.
    \choice Their score is 80.
\end{choices}

\question[1] If every value in a data set is multiplied by 3, then the standard deviation:
\begin{choices}
    \choice Stays the same
    \choice Increases by 3
    \correctchoice Is multiplied by 3
    \choice Is multiplied by 9
\end{choices}

\question[1] If a constant $k$ is added to every data value, then the mean:
\begin{choices}
    \choice Stays the same
    \correctchoice Increases by $k$
    \choice Is multiplied by $k$
    \choice Decreases by $k$
\end{choices}

\question[1] A box-and-whisker plot displays:
\begin{choices}
    \choice Mean, median, and mode only
    \correctchoice Minimum, $Q_1$, median, $Q_3$, and maximum
    \choice Only outliers
    \choice Individual data values
\end{choices}

\question[1] Which of the following represents the most spread out data set?
\begin{choices}
    \choice $\sigma = 0.5$
    \choice $\sigma = 1.2$
    \correctchoice $\sigma = 4.8$
    \choice $\sigma = 2.1$
\end{choices}

\question[1] If a data set has a standard deviation of 0, then:
\begin{choices}
    \choice The mean is 0
    \correctchoice All values in the data set are equal
    \choice The data set has only one value
    \choice The data set is symmetric
\end{choices}

\question[1] The range of a data set is:
\begin{choices}
    \choice The most frequently occurring value
    \correctchoice The difference between the maximum and minimum values
    \choice $Q_3 - Q_1$
    \choice The middle value when sorted
\end{choices}


\section*{Part B --- Short Answer (33 marks)}

\question The following data represents the number of goals scored per game by a hockey team
over 12 games:
\[3,\; 1,\; 4,\; 2,\; 6,\; 3,\; 5,\; 3,\; 2,\; 8,\; 1,\; 4\]
\begin{parts}
    \part[2] Find the mean, median, and mode.
    \part[2] Find $Q_1$, $Q_3$, and the IQR.
    \part[3] Calculate the population standard deviation. Show all steps.
    \part[2] Using the $1.5 \times \text{IQR}$ rule, identify any outliers.
\end{parts}
\begin{solution}[11cm]
Sorted data: $1, 1, 2, 2, 3, 3, 3, 4, 4, 5, 6, 8$

\textbf{(a)}
$\bar{x} = \dfrac{1+1+2+2+3+3+3+4+4+5+6+8}{12} = \dfrac{42}{12} = \mathbf{3.5}$

Median: average of 6th and 7th values $= \dfrac{3+3}{2} = \mathbf{3}$\quad Mode: $\mathbf{3}$

\textbf{(b)} Lower half: $1,1,2,2,3,3$ $\Rightarrow$ $Q_1 = \dfrac{2+2}{2} = \mathbf{2}$

Upper half: $3,4,4,5,6,8$ $\Rightarrow$ $Q_3 = \dfrac{4+5}{2} = \mathbf{4.5}$

$\text{IQR} = 4.5 - 2 = \mathbf{2.5}$

\textbf{(c)} $\mu = 3.5$.
\[
\sigma^2 = \frac{1}{12}\sum(x_i - 3.5)^2
= \frac{(1-3.5)^2\times2 + (2-3.5)^2\times2 + (3-3.5)^2\times3 + (4-3.5)^2\times2 + (5-3.5)^2 + (6-3.5)^2 + (8-3.5)^2}{12}
\]
\[
= \frac{12.5 + 4.5 + 0.75 + 0.5 + 2.25 + 6.25 + 20.25}{12}
= \frac{47}{12} \approx 3.917
\]
\[
\sigma \approx \mathbf{1.979}
\]

\textbf{(d)} Lower fence: $Q_1 - 1.5\,\text{IQR} = 2 - 3.75 = -1.75$.
Upper fence: $Q_3 + 1.5\,\text{IQR} = 4.5 + 3.75 = 8.25$.\\
All data values lie within $[-1.75,\; 8.25]$, so there are \textbf{no outliers}.
\end{solution}

\question Two students each wrote 8 assignments. Their scores are:

\begin{center}
\begin{tabular}{lcccccccc}
\toprule
\textbf{Student A} & 55 & 62 & 68 & 70 & 72 & 75 & 80 & 82 \\
\textbf{Student B} & 40 & 58 & 65 & 72 & 74 & 78 & 85 & 92 \\
\bottomrule
\end{tabular}
\end{center}

\begin{parts}
    \part[2] Calculate the mean for each student.
    \part[3] Calculate the population standard deviation for each student.
    \part[3] Compare the two students' performances using mean and standard deviation.
    Who is more consistent? Who performed better overall?
\end{parts}
\begin{solution}[9cm]
\textbf{(a)} $\bar{A} = \dfrac{55+62+68+70+72+75+80+82}{8} = \dfrac{564}{8} = \mathbf{70.5}$

$\bar{B} = \dfrac{40+58+65+72+74+78+85+92}{8} = \dfrac{564}{8} = \mathbf{70.5}$

\textbf{(b)} For A: deviations from 70.5: $-15.5, -8.5, -2.5, -0.5, 1.5, 4.5, 9.5, 11.5$
\[
\sigma_A^2 = \frac{240.25 + 72.25 + 6.25 + 0.25 + 2.25 + 20.25 + 90.25 + 132.25}{8}
= \frac{564}{8} = 70.5
\]
$\sigma_A = \sqrt{70.5} \approx \mathbf{8.40}$

For B: deviations from 70.5: $-30.5, -12.5, -5.5, 1.5, 3.5, 7.5, 14.5, 21.5$
\[
\sigma_B^2 = \frac{930.25+156.25+30.25+2.25+12.25+56.25+210.25+462.25}{8}
= \frac{1860}{8} = 232.5
\]
$\sigma_B = \sqrt{232.5} \approx \mathbf{15.25}$

\textbf{(c)} Both students have the same mean ($\bar{x} = 70.5$), so their overall performance is
equal. However, Student A is more consistent ($\sigma_A \approx 8.40$ vs $\sigma_B \approx 15.25$),
meaning Student B's scores are much more variable.
\end{solution}

\question The ages of patients in a clinic on a given day are:
\[22,\; 25,\; 31,\; 35,\; 38,\; 42,\; 45,\; 48,\; 52,\; 67\]
\begin{parts}
    \part[2] Find the five-number summary: Min, $Q_1$, Median, $Q_3$, Max.
    \part[2] Sketch a clearly labelled box-and-whisker plot.
    \part[2] A new patient aged 85 arrives. Recalculate the five-number summary and
    state whether 85 is an outlier.
\end{parts}
\begin{solution}[9cm]
Sorted (already sorted). $n = 10$.

\textbf{(a)} Min $= 22$, Max $= 67$.\\
Median $= \dfrac{38+42}{2} = 40$.\\
$Q_1 = \dfrac{25+31}{2} = 28$. \quad $Q_3 = \dfrac{48+52}{2} = 50$.\\
\textbf{Five-number summary: $\{22,\; 28,\; 40,\; 50,\; 67\}$}

\textbf{(b)} Box from 28 to 50, median line at 40, whiskers to 22 and 67.

\textbf{(c)} $n=11$: Median $= 42$. Lower 5: $22,25,31,35,38\Rightarrow Q_1=31$.
Upper 5: $45,48,52,67,85\Rightarrow Q_3=52$. IQR $= 21$.\\
Upper fence: $52 + 1.5(21) = 52 + 31.5 = 83.5$.\\
Since $85 > 83.5$, \textbf{85 is an outlier}.
\end{solution}

\question A teacher records the following test scores (out of 100) for her class of 10 students:
\[55,\; 60,\; 65,\; 70,\; 72,\; 75,\; 80,\; 85,\; 90,\; 98\]
\begin{parts}
    \part[2] What score corresponds to the 70th percentile?
    \part[2] What is the percentile rank of a student who scored 72?
    \part[2] The teacher decides to curve all marks by multiplying by 1.1.
    What are the new mean and standard deviation?
\end{parts}
\begin{solution}[7cm]
\textbf{(a)} 70th percentile: position $= 0.70 \times 10 = 7$. Since this is a whole number,
average the 7th and 8th values: $\dfrac{80+85}{2} = \mathbf{82.5}$

\textbf{(b)} Scores at or below 72: $\{55, 60, 65, 70, 72\}$ = 5 values.
\[
\text{Percentile rank} = \frac{5}{10} \times 100 = \mathbf{50\text{th percentile}}
\]

\textbf{(c)} Original mean: $\bar{x} = \dfrac{750}{10} = 75$.\\
Original $\sigma = \sqrt{\dfrac{\sum(x_i-75)^2}{10}}$. Deviations: $-20,-15,-10,-5,-3,0,5,10,15,23$.
$\sum d^2 = 400+225+100+25+9+0+25+100+225+529 = 1638$.
$\sigma = \sqrt{163.8} \approx 12.80$.

After multiplying by 1.1: new mean $= 75 \times 1.1 = \mathbf{82.5}$;
new $\sigma = 12.80 \times 1.1 \approx \mathbf{14.08}$.
\end{solution}

\end{questions}
\end{document}
